%=============================================================================!
% 12.5  THE DISTRIBUTION OF VELOCITIES                                        !
%=============================================================================!
\section{The distribution of velocities}
\label{sec:sm-maxwell}

\Cref{sec:md-start} drew starting velocities from the normal distribution
and promised a reason. This section derives the distribution of the
velocities of atoms at a temperature $T$ from the Boltzmann factor, then
the distributions of their speeds and kinetic energies, and measures all
three in liquid argon.

\subsection*{Each component is Gaussian}

We first take an unconstrained canonical system, with no condition on its
total momentum. The energy of $N$ atoms is $H = \sum_{i,x}p_{ix}^2/2m_i + U(\vb{r}^N)$,
where the sum runs over the atoms $i$ and, with $x$ standing for each in
turn, the three directions $x$, $y$ and $z$. The
exponential of a sum is a product, so the Boltzmann factor of
\cref{eq:sm-boltzmann} splits into one factor for the positions and one
for each component of momentum:
\begin{equation*}
   \mathrm{e}^{-\beta H} = \mathrm{e}^{-\beta U(\vb{r}^N)}\prod_{i,x}
   \mathrm{e}^{-\beta p_{ix}^2/2m_i} .
\end{equation*}
A probability density that is a product of factors, each depending on
one quantity alone, describes quantities that are independent, so every
component of every momentum is independent of the others and of the
positions. With $p_{ix} = m_iv_{ix}$, the density of one component of
velocity is proportional to $\mathrm{e}^{-m_iv_{ix}^2/2\kB T}$, which is
the Gaussian of \cref{eq:sm-gauss} with mean zero and variance
\begin{equation}
   \sigma_v^2 = \frac{\kB T}{m_i} .
   \label{eq:sm-maxwell}
\end{equation}
This is the Maxwell-Boltzmann distribution of velocities. Two things in it
matter for simulation. It holds whatever the forces: the velocities of
atoms in a gas, a liquid or a crystal at the same temperature have the
same distribution, in classical mechanics. And the spread depends on the
mass: a lithium atom at \SI{300}{\kelvin} has $\sigma_v = \sqrt{\kB T/m} =
\sqrt{0.025852/(6.94\times103.6427)} =
\SI{0.005995}{\angstrom\per\femto\second}$, the factor 103.6427 of
\cref{eq:en-mv2} turning \si{\amu\angstrom\squared\per\femto\second\squared}
into \si{\electronvolt}, and an argon atom
\SI{0.002499}{\angstrom\per\femto\second}, smaller by $\sqrt{39.948/6.94}
= 2.399$, the ratio of the square roots of their masses. Dividing the normal numbers by
$\sqrt{m_i}$, as \texttt{starting\_velocities} of \cref{sec:md-start}
did, gives each atom this spread up to a common factor, which the
scaling to the chosen energy then fixed.

\texttt{mdlab.statmech.thermal\_velocities} draws velocities at a
temperature directly:

\begin{code}
    m = np.asarray(masses, dtype=float)
    xi = rng.standard_normal((len(m), 3))
    v = xi * np.sqrt(kb * temperature / (m * mv2_to_energy))[:, None]
    if remove_drift:
        v = v - (m @ v) / m.sum()
    return v
\end{code}

\noindent Line 2 draws $3N$ numbers of mean 0 and variance 1, and line 3
multiplies each atom's by its $\sigma_v$, with \texttt{mv2\_to\_energy}
the factor of \cref{eq:en-mv2} that turns $mv^2$ into
\si{\electronvolt}. Line 5 removes the drift of the centre of mass, as in
\cref{sec:md-start}. ASE's \texttt{thermalize\_momenta} draws the numbers
in the same order, so the same starting number of the random generator,
its seed, gives the same velocities: for 50
atoms of random masses at \SI{300}{\kelvin} they differ by at most
\num{1.7e-7} of their size. That is half the relative difference of
\num{3.4e-7} between the two libraries' values of $\kB$, ASE's being the
value recommended for the constants in 2014, since $\sigma_v$ goes as
$\sqrt{\kB}$ and a small relative change in a quantity changes its square
root by half as much.

\subsection*{Speeds and kinetic energies}

The speed $v = \lvert\vb{v}\rvert$ of an atom combines three independent
components. The probability that the velocity lies in a small box $\dd
v_x\,\dd v_y\,\dd v_z$ is the product of the three Gaussians, each with the
factor $1/(\sigma_v\sqrt{2\pi}) = (m/2\pi\kB T)^{1/2}$, times the box, and
with $v_x^2 + v_y^2 + v_z^2 = v^2$ it is
\begin{equation*}
   \Bigl(\frac{m}{2\pi\kB T}\Bigr)^{3/2}\mathrm{e}^{-mv^2/2\kB T}\,
   \dd v_x\,\dd v_y\,\dd v_z ,
\end{equation*}
which depends on the velocity only through its size. The velocities with
speeds between $v$ and $v + \dd v$ fill a spherical shell of area $4\pi
v^2$ and thickness $\dd v$, so the density of the speed is
\begin{equation}
   \mathcal{P}(v) = 4\pi v^2\Bigl(\frac{m}{2\pi\kB T}\Bigr)^{3/2}
   \mathrm{e}^{-mv^2/2\kB T} .
   \label{eq:sm-speed}
\end{equation}
The factor $v^2$ makes it zero at $v = 0$, since few velocities have all
three components near zero. It peaks at $\sqrt2\,\sigma_v$ and has the
mean $\sqrt{8/\pi}\,\sigma_v$ (\cref{ex:sm-speeds}); the mean speed of
lithium at \SI{300}{\kelvin} is \SI{0.00957}{\angstrom\per\femto\second},
or \SI{957}{\metre\per\second}, with \SI{1}{\angstrom\per\femto\second}
equal to \SI{e5}{\metre\per\second}, and that of argon
\SI{399}{\metre\per\second}.

An atom's kinetic energy $K_i = \frac12mv^2$ follows by a change of
variable. Between $K_i$ and $K_i + \dd K_i$ lie the speeds between $v$
and $v + \dd v$ with $\dd K_i = mv\,\dd v$, and the two hold the same
probability, $\mathcal{P}(K_i)\,\dd K_i = \mathcal{P}(v)\,\dd v$, so
$\mathcal{P}(K_i) = \mathcal{P}(v)/(mv)$. With $mv^2/2\kB T = K_i/\kB T$
and $v\sqrt m = \sqrt{2K_i}$,
\begin{equation*}
   \frac{\mathcal{P}(v)}{mv} = \frac{4\pi v\sqrt m}{(2\pi\kB T)^{3/2}}\,
   \mathrm{e}^{-K_i/\kB T}
   = \frac{4\pi\sqrt2}{(2\pi)^{3/2}}\,(\kB T)^{-3/2}\sqrt{K_i}\,
   \mathrm{e}^{-K_i/\kB T} ,
\end{equation*}
and $4\pi\sqrt2/(2\pi)^{3/2} = 4\pi\sqrt2/(2\sqrt2\,\pi^{3/2}) =
2/\sqrt\pi$, so
\begin{equation}
   \mathcal{P}(K_i) = \frac{2}{\sqrt\pi}\,(\kB T)^{-3/2}\sqrt{K_i}\,
   \mathrm{e}^{-K_i/\kB T} ,
   \label{eq:sm-atom-energy}
\end{equation}
the same for every mass. Its mean is $\frac32\kB T$ and its variance
$\frac32(\kB T)^2$, as \cref{sec:sm-kinetic} shows for any number of
freedoms.

\subsection*{Measured in a liquid}

\Cref{fig:sm-maxwell} tests the three densities on the liquid argon of
\cref{fig:sm-ergodic}(b), 256 atoms over \SI{50}{\pico\second}, with the
velocities saved every \SI{50}{\femto\second}: 768\,768 components from
1001 frames. The run is at fixed energy, not in contact with a bath, but
each atom has the other 255 around it, which act as its bath. The
kinetic temperature of the run, defined in \cref{sec:sm-equipartition},
is \SI{134.63}{\kelvin}, at which \cref{eq:sm-maxwell} gives $\sigma_v =
\SI{0.001674}{\angstrom\per\femto\second}$; the components have the
standard deviation \SI{0.001671}{\angstrom\per\femto\second}, and
\cref{sec:sm-equipartition} accounts for the difference. The mean speed is
\SI{0.002668}{\angstrom\per\femto\second} against 0.002671 from
\cref{eq:sm-speed}, and the kinetic energies of single atoms have the
mean $1.494\,\kB T$ and the variance $1.479\,(\kB T)^2$, against
$\frac32$ for both.

\begin{figure}[htb]
\centering
\includegraphics{ch12_ensembles/maxwell}
\caption{Velocities in liquid argon at \SI{134.6}{\kelvin}: histograms of
256 atoms over \SI{50}{\pico\second}, with velocities in
\si{\angstrom\per\pico\second}, a thousand times their values in
\si{\angstrom\per\femto\second}. (a) Every component of velocity, against
the Gaussian of variance $\kB T/m$ (dashed). (b) The speeds, against
\cref{eq:sm-speed}. (c) The kinetic energies of single atoms, in units of
$\kB T$, against \cref{eq:sm-atom-energy}, a density per unit of $\kB
T$.}
\label{fig:sm-maxwell}
\end{figure}

A Gaussian has a property that the histogram shows only roughly and a
single number shows well. The fourth moment of the Gaussian of mean zero
is $\ens{v^4} = 3\sigma_v^4$ (\cref{ex:sm-fourth}), so the excess
kurtosis
\begin{equation*}
   \frac{\ens{u^4}}{\ens{u^2}^2} - 3 ,
\end{equation*}
computed over the mass-weighted components $u = \sqrt{m}\,v$, which share
one variance $\kB T$ whatever the masses, is zero for velocities with the
Maxwell-Boltzmann distribution; pooled components of different variances
would give more than zero even if each were Gaussian. Components spread
evenly over an interval give $-1.2$ (\cref{ex:sm-fourth}). Over the
768\,768 components of the liquid it is $-0.0072$, small against the
$-1.2$ of an even spread. \texttt{mdlab.diagnostics.excess\_kurtosis}
computes it, and \cref{sec:sm-trajectory} uses it to catch velocities that
are not yet Maxwellian.

\begin{keyidea}
In the unconstrained canonical ensemble each component of each velocity
is independent and Gaussian, with mean zero and variance $\kB T/m$, whatever the forces.
The speed has the density $4\pi v^2(m/2\pi\kB T)^{3/2}\mathrm{e}^{-mv^2/
2\kB T}$, and an atom's kinetic energy a density with mean $\frac32\kB T$
for every mass. The excess kurtosis of the mass-weighted components is
zero for Maxwellian velocities. Removing the total momentum or holding
bond lengths constrains these components, as the next section explains.
\end{keyidea}

\notebook{12}{draw velocities at any temperature and mass, and check them against ASE}

\begin{selfcheck}
What is $\sigma_v$ for a lithium atom at \SI{600}{\kelvin}?
\end{selfcheck}
